%% main.tex — FIXTURE A of the underclaiming/buried-lede regression pair.
%%
%% Issue #1046 / #1048. This paper is deliberately RIGOROUS AND
%% UNDERCLAIMING: its method, numbers, ablations, seeds, tables, threats
%% section, and bibliography are IDENTICAL to Fixture B's
%% (../../../bold-synthesis-labeled/...), but the organizing idea is
%% demoted to the last paragraph of the Discussion and never appears in
%% the abstract, the title, or the contribution list.
%%
%% Everything here is SYNTHETIC — the organization, the traces, the
%% measurements, and every bibliography entry are invented for the
%% fixture. Do not cite any of it.
%%
%% The Method-through-Experiments span is byte-identical to Fixture B's
%% and is pinned by tests/test_paper_underclaiming_fixtures.py, which is
%% what makes the score divergence attributable to framing alone. Edit
%% both fixtures or neither.

\documentclass{anvil-paper}

\title{An Empirical Characterization of Remote Cache Miss Behavior in Distributed Build Systems}
\author{Author Names Withheld}
\date{\today}

\begin{document}
\maketitle

\begin{abstract}
Distributed build systems rely on content-addressed remote caches to avoid
recomputing unchanged targets, and prior deployment reports place aggregate
hit rates above ninety percent. Less is known about how the residual misses
are distributed across a repository. We report a fourteen-week measurement
study of 2.1 million build invocations across three production repositories,
in which misses are found to be unevenly distributed: 8.3\% of source files
account for 71.4\% of remote-cache misses. We additionally implement a
prototype pre-warming heuristic that ranks files by version-control churn and
co-change centrality, and observe a 22.4\% median reduction in end-to-end
build latency in one eight-week deployment. The characterization is
descriptive, the prototype is preliminary, and we do not claim that either
result generalizes beyond the environment measured here; threats to validity
are discussed in Section~\ref{sec:discussion}.
\end{abstract}

\section{Introduction}
\label{sec:intro}

Build systems that distribute compilation across a shared cluster have become
standard at organizations above a few hundred engineers~\cite{okonkwo2019remote}.
Such systems key a content-addressed cache on the transitive hash of each
action's inputs, so an action whose inputs are unchanged is served from cache
rather than recomputed. Reported aggregate hit rates are high, and the
literature has accordingly concentrated on cache capacity, eviction, and
transport rather than on the composition of the miss set.

Measurement of production build clusters is difficult for reasons that
constrain what can responsibly be claimed from it. Traces are proprietary;
workloads are entangled with organizational practices such as branch cadence
and review latency; and the metric developers care about, wall-clock time to a
usable artifact, depends on cluster contention that varies hour to hour. Prior
measurement work has been careful to scope its conclusions to the
organizations studied~\cite{vargas2018buildlatency}, and we adopt the same
posture here.

We are likewise constrained. Our three repositories belong to a single
organization, share a single CI configuration lineage, and use one build tool.
The deployment arm of the study ran on the same cluster that produced the
observational trace, so the two are not independent samples. We report what we
measured, we report the ablations that isolate which signal contributes what,
and we resist extrapolating to build environments we did not observe.

Within those limits, this paper makes three contributions:

\begin{itemize}
  \item A fourteen-week trace of 2.1 million build invocations across three
        production repositories, released in anonymized form.
  \item A characterization of the miss set showing that misses concentrate in
        a small minority of source files, and that membership in that minority
        is predictable from version-control history with AUC 0.87.
  \item A prototype pre-warming heuristic and a preliminary eight-week
        deployment evaluation, with an ablation separating the churn and
        centrality signals.
\end{itemize}

Section~\ref{sec:related} positions the work against prior caching and
change-impact literature, Section~\ref{sec:method} describes the trace and the
predictor, Section~\ref{sec:experiments} reports the measurements, and
Section~\ref{sec:discussion} discusses limitations.

\section{Related Work}
\label{sec:related}

\paragraph{Build caching and remote execution.}
Content-addressed caching for build actions is well established in practice
and in the literature~\cite{okonkwo2019remote}, as are the transport and
eviction questions that follow from it. Our measurement adds a distributional
view of the miss set to that body of work but does not modify the caching
mechanism itself.

\paragraph{Prefetching in continuous integration.}
Lindqvist and Osei~\cite{lindqvist2021prefetch} prefetch CI artifacts by
recency, warming whatever was built most recently. Their recency signal is the
baseline we compare against in Section~\ref{sec:experiments}. Our heuristic
differs only in the ranking function it uses to choose targets; the warming
mechanism is theirs.

\paragraph{Co-change mining.}
The churn and co-change measures we use are taken directly from the mining
literature. Ferreira and Adeyemi~\cite{ferreira2016cochange} define the
co-change graph and its impact-analysis use; Nakamura et
al.~\cite{nakamura2020churn} establish the churn windows we adopt without
modification; and the centrality measure is the standard eigenvector
formulation of Imani and Pereira~\cite{imani2017graphcentrality}. Our ranking
function is a straightforward re-application of these measures to a new input,
and we claim no novelty in the predictor itself.

\paragraph{History-aware test selection.}
Bhattacharya and Lund~\cite{bhattacharya2022testsel} select regression tests
using historical failure data. That work targets test selection rather than
cache warming, so we do not compare against it directly.

\section{Method}
\label{sec:method}

\subsection{Trace collection}

Each build invocation is a directed acyclic graph of actions; the build tool
emits one record per action with the action's cache key, the resolved input
hashes, the outcome (\texttt{hit}, \texttt{miss}, or \texttt{local}), and the
wall-clock duration. We collected every such record for fourteen consecutive
weeks across three repositories: R1, a 41k-file monorepo; R2, a 9.8k-file
service repository; and R3, a 6.2k-file embedded firmware repository. The
trace contains 2.1 million build invocations comprising 903 million action
records.

A miss is attributed to a source file when that file's content hash appears in
the action's resolved input set and differs from the hash present at the last
observed hit for the same action identity. Actions whose input set changed in
more than one source file are attributed fractionally, splitting the miss
evenly across the changed inputs; Section~\ref{sec:discussion} reports the
sensitivity of our results to that choice.

\subsection{Predictor}

For each source file and each day we compute four features: the count of
commits touching the file in the preceding 30 days; the eigenvector centrality
of the file in the 90-day co-change graph~\cite{imani2017graphcentrality};
the reverse-dependency fan-out of the file in the build graph; and the mean
diff size across the file's commits in the preceding 30 days. A logistic
regression predicts whether the file will appear in the next day's
top-decile miss set. We fit with L2 regularization ($\lambda = 0.01$), five-fold
cross-validation grouped by week to avoid temporal leakage, and a fixed seed
(20260214) recorded in the released artifact.

\subsection{Pre-warming}

The pre-warmer ranks files by predicted probability, takes the top $k$, and
speculatively builds their downstream targets during idle cluster capacity.
The parameter $k$ is set by a controller that holds speculative work within a
6\% cluster-CPU budget; the controller and its budget are described in the
released configuration. No developer-visible behavior changes: speculative
work is preemptible and is discarded when a real build arrives.

\section{Experiments}
\label{sec:experiments}

\subsection{Miss concentration}

Table~\ref{tab:concentration} reports the concentration of the miss set. Across
the pooled trace, 8.3\% of source files account for 71.4\% of remote-cache
misses. Concentration is present in all three repositories and is strongest in
the monorepo R1, where 6.1\% of files account for 74.8\% of misses.

\begin{table}[t]
  \centering
  \begin{tabular}{lrrr}
    \toprule
    Repository & Files & Miss-heavy files (\%) & Misses covered (\%) \\
    \midrule
    R1 (monorepo)  & 41{,}200 & 6.1 & 74.8 \\
    R2 (service)   & 9{,}800  & 9.7 & 68.2 \\
    R3 (firmware)  & 6{,}200  & 11.4 & 66.9 \\
    \midrule
    Pooled         & 57{,}200 & 8.3 & 71.4 \\
    \bottomrule
  \end{tabular}
  \caption{Concentration of remote-cache misses. ``Miss-heavy files'' is the
  smallest set of files, ranked by attributed misses, covering the reported
  share of all misses over the fourteen-week trace.}
  \label{tab:concentration}
\end{table}

\subsection{Predictability}

The logistic regression attains AUC 0.87 (95\% CI [0.86, 0.88], five-fold
grouped cross-validation) on next-day top-decile miss-set membership. The
recency-only baseline of Lindqvist and Osei~\cite{lindqvist2021prefetch}
attains AUC 0.71 on the same folds. Table~\ref{tab:ablation} reports the
feature ablation.

\begin{table}[t]
  \centering
  \begin{tabular}{lrr}
    \toprule
    Signal & AUC & Median latency reduction (\%) \\
    \midrule
    Recency only (baseline)      & 0.71 & 5.3 \\
    Churn only                   & 0.79 & 9.1 \\
    Centrality only              & 0.76 & 7.4 \\
    Churn $+$ centrality (full)  & 0.87 & 22.4 \\
    \bottomrule
  \end{tabular}
  \caption{Feature ablation. Latency reductions are medians over the
  deployment arm described in Section~\ref{sec:experiments}.}
  \label{tab:ablation}
\end{table}

\subsection{Deployment}

We deployed the pre-warmer for eight weeks, cluster-randomized by engineer
across 1{,}412 engineers and 612{,}000 builds. Median end-to-end build latency
fell 22.4\% (95\% CI [20.9, 23.8]) in the treatment arm; p90 latency fell
31.1\%. Cluster CPU rose 5.8\%, within the configured budget. No regression
was observed in build correctness auditing, which compares speculative outputs
against the subsequently computed real outputs on a 1\% sample.

The combined signal outperforms the sum of its parts: churn alone yields 9.1\%
and centrality alone 7.4\%, while the combination yields 22.4\%. The two
signals select different files. Churn ranks recently edited leaves highly;
centrality ranks stable files whose edits propagate widely. Neither ordering
alone covers the miss set that drives tail latency.

\section{Discussion}
\label{sec:discussion}

\paragraph{Threats to validity.}
The three repositories belong to one organization and share a CI lineage, so
practices such as branch cadence are held constant rather than varied. The
deployment arm ran on the cluster that produced the observational trace, so
the two are not independent. Fractional miss attribution for multi-input
changes is a modeling choice; re-running the analysis with
first-changed-input attribution moves pooled coverage from 71.4\% to 68.9\%
and does not change the ordering in Table~\ref{tab:ablation}. Finally, the
deployment measured wall-clock latency under the cluster's ordinary contention
and not under saturation, where speculative work is likelier to displace real
work.

\paragraph{Cost.}
The 5.8\% CPU overhead is a real cost, and organizations whose clusters run
near saturation should expect the trade to be less favorable than it was here.

\paragraph{A possible reading.}
More speculatively, one might read these measurements as suggesting that
version-control history is an input to build \emph{scheduling} rather than
only to code analysis, and that a build system might treat repository history
as a first-class scheduling signal alongside the dependency graph. The
ablation is at least consistent with such a reading, since the two history
signals contribute non-redundantly. We do not develop that framing here and
leave it to future work.

\section{Conclusion}
\label{sec:conclusion}

We characterized the distribution of remote-cache misses in a fourteen-week
production trace and found that misses concentrate in a small minority of
source files that is predictable from version-control history. A prototype
pre-warmer built on that predictor reduced median build latency by 22.4\% in
one eight-week deployment, at a 5.8\% cluster-CPU cost. The characterization
is descriptive and the deployment is preliminary.

\bibliographystyle{plainnat}
\bibliography{refs}

\end{document}
